\section{关键算法与模型原理}

本项目的核心 NLP 任务包括文本语言识别、情感三分类、课程维度识别、关键词提取和结构化建议生成。
其中，情感分类承担主要可量化评估任务；
课程维度和证据识别保证结果可解释；
LLM 只负责把结构化结果组织成自然语言建议。

\subsection{双语预处理与情感文本标准化}

基础预处理先删除 URL、异常符号和多余空白，再识别文本语言。
中文使用 jieba 分词；
英文使用正则表达式分词并过滤常见停用词。
为了保留课程领域信息，系统不会删除“作业、考试、实验、讲解、收获”等对教学分析有意义的词。

真实评论中经常出现 \code{tbh}、\code{idk}、\code{cant}、\code{crs} 等缩写；
也会出现 \code{excelnt}、\code{confuzing}、\code{uncler} 等拼写错误。
系统在情感分类前增加独立标准化层，将常见缩写展开；
同时只对情感词执行受限纠错。
自动纠错必须满足候选唯一、相同词首，且最多只有一次漏字、多字或相邻字母换位。
合法词形和技术名词受到保护。
部分转换示例如表~\ref{tab:normalization-examples}。

\begin{table}[H]
  \centering
  \caption{噪声情感文本标准化示例}
  \label{tab:normalization-examples}
  \small
  \begin{tabularx}{0.92\textwidth}{Y Y}
    \toprule
    原始表达 & 送入情感分类器的标准化表达 \\
    \midrule
    \texttt{Tbh this crs is dificult} & \texttt{To be honest, this course is difficult} \\
    \texttt{I cant recomend this crs} & \texttt{I cannot recommend this course} \\
    \texttt{borng and uncler} & \texttt{boring and unclear} \\
    \texttt{goooood} & \texttt{good} \\
    \bottomrule
  \end{tabularx}
\end{table}

标准化后的文本只用于情感分类。
原始文本继续用于页面展示、关键词、课程维度和证据提取。
这种“双文本通道”避免了自动纠错修改学生原话；
也保证最终证据能够与用户输入直接对应。

\subsection{multilingual BERT 情感分类}

BERT 使用 Transformer 编码器；
通过双向自注意力建模词语在上下文中的关系
\citep{devlin2019bert,vaswani2017attention}。
对于输入 Token 序列 \(x_1,x_2,\ldots,x_n\)，模型得到各位置的上下文表示；
并使用 \code{[CLS]} 位置的最终表示 \(h_{\mathrm{CLS}}\) 完成三分类：

\begin{equation}
  \boldsymbol{p}
  = \operatorname{softmax}\left(W h_{\mathrm{CLS}} + \boldsymbol{b}\right),
\end{equation}

其中 \(\boldsymbol{p}\) 分别表示消极、中性和积极三个类别的概率。
最终标签为概率最大的类别。
项目采用 \code{bert-base-multilingual-cased}，利用其多语言词表同时处理中文和英文。

\begin{table}[H]
  \centering
  \caption{BERT 训练配置}
  \label{tab:bert-training}
  \small
  \begin{tabularx}{\textwidth}{L{3.2cm} Y L{3.2cm} Y}
    \toprule
    参数 & 数值 & 参数 & 数值 \\
    \midrule
    训练/验证/测试集 & 5472/1368/2280 &
    随机种子 & 42 \\
    训练轮数 & 2 &
    批大小 & 8 \\
    最大序列长度 & 160 &
    学习率 & \(2\times10^{-5}\) \\
    权重衰减 & 0.01 &
    最优模型指标 & Macro-F1 \\
    \bottomrule
  \end{tabularx}
\end{table}

验证集用于每轮训练后的模型选择。
独立测试集只在训练完成后评估一次。
训练时若 CUDA 可用则启用 GPU；
若不可用，也可以在 CPU 上运行。
部署阶段采用懒加载方式，第一次分析时加载模型，后续请求复用同一实例。

\subsection{长文本滑动窗口推理}

直接截断长评价会丢失后半部分信息。
系统将超过最大长度的 Token 序列按滑动窗口切分。
默认最大长度为 160，重叠 32 个 Token，单条评价最多保留 16 个窗口。
整条评价的加权分类概率为：

\begin{equation}
  \label{eq:weighted-window}
  \boldsymbol{p}
  = \frac{\sum_{i=1}^{m} w_i \boldsymbol{p}_i}
  {\sum_{i=1}^{m} w_i}.
\end{equation}

其中，\(\boldsymbol{p}_i\) 表示第 \(i\) 个窗口的分类概率；
\(w_i\) 表示该窗口的有效 Token 数；
\(m\) 表示窗口数量。
使用有效 Token 数加权，可以避免短尾窗口和完整窗口具有相同影响。
后端会记录窗口数量、Token 数量和是否达到窗口上限。

\subsection{规则校正与分层回退}

BERT 能理解大部分上下文；
但对课程评价中的混合态度、局部否定和固定业务定义仍可能产生偏差。
规则模块统计中英文正负信号，并检测“但是、但、不过、but、however”等转折结构。
其主要处理逻辑包括：

\begin{itemize}
  \item 正负信号同时出现且相对平衡时，将结果校正为中性；
  \item “不差、并非没有帮助”等否定消极表达作为弱积极信号；
  \item “不清楚、not useful”等否定积极表达作为消极信号；
  \item 建议型表达默认保持中性，避免把“希望增加代码示例”误判为强消极；
  \item 对部分明显讽刺或隐含抱怨使用有限语义规则，但不声称完整解决讽刺识别。
\end{itemize}

系统默认采用 \code{AUTO} 后端，执行顺序为：

\begin{equation}
  \text{BERT + 规则校正}
  \rightarrow \text{TF-IDF Logistic Regression}
  \rightarrow \text{纯规则模型}.
\end{equation}

该顺序将语义效果较好的 BERT 放在首位；
同时保留体积小、启动快的传统模型，以及无需模型文件的规则兜底。

\subsection{传统情感模型}

传统模型将预处理后的文本转换为 TF-IDF 向量。
词 \(t\) 在文档 \(d\) 中的权重表示为：

\begin{equation}
  \operatorname{tfidf}(t,d)
  = \operatorname{tf}(t,d)
  \log\frac{N}{\operatorname{df}(t)+1},
\end{equation}

其中 \(N\) 为文档总数，\(\operatorname{df}(t)\) 为包含词 \(t\) 的文档数。
项目使用一元与二元词组、最多 30000 个特征；
并比较 Dummy、Logistic Regression、Naive Bayes 和 Linear SVM。
最佳传统模型为 Logistic Regression。

\begin{table}[H]
  \centering
  \caption{传统模型与 BERT 测试结果}
  \label{tab:model-results}
  \small
  \begin{tabular}{lcc}
    \toprule
    模型 & Accuracy & Macro-F1 \\
    \midrule
    Dummy Most Frequent & 0.3333 & 0.1667 \\
    Naive Bayes & 0.6772 & 0.6785 \\
    Linear SVM & 0.6811 & 0.6795 \\
    Logistic Regression & 0.6904 & 0.6906 \\
    multilingual BERT & \textbf{0.7461} & \textbf{0.7457} \\
    \bottomrule
  \end{tabular}
\end{table}

BERT 相比 Logistic Regression 的 Accuracy 提高约 5.57 个百分点。
Macro-F1 提高约 5.51 个百分点。
BERT 对积极类效果最好；
中性类由于混合评价和建议型表达较多，仍是最难识别的类别。

\subsection{课程维度与证据识别}

系统使用面向课程场景的双语领域词典识别六个维度：
教学内容、授课方式、作业任务、考试安排、实验实践和学习收获。
例如出现“assignments”和“deadline”时，系统识别为“作业任务”；
出现“setup”和“code”时，系统识别为“实验实践”。

对于每个命中维度，系统不仅返回维度名称；
还记录最先出现的关键词、出现位置和关键词附近的原文片段。
维度结果按首次出现位置排序，因此能够保留评价的叙述顺序。
这种方法没有采用复杂主题模型；
但在固定课程场景中实现简单、结果稳定，也便于人工检查和扩展词典。

\subsection{关键词提取与证据输出}

关键词模块统计预处理后词项频次并过滤停用词。
在单条分析页面，关键词用于概括当前评价中的核心表达；
在批量分析页面，关键词用于定位待关注评论中的高频问题。

关键词、课程维度和原文证据共同进入结构化结果。
它们的作用不是替代教师阅读原文，而是帮助教师更快定位“学生在说什么”和“问题属于哪个课程环节”。

\subsection{结构化建议生成}

大语言模型不直接接收孤立评论并自由判断。
系统先生成情感倾向、课程维度、关键词和证据片段；
再把这些结构化信息送入建议生成模块。
提示词要求输出 summary、problems、suggestions 和 risk level 等固定字段。
温度设置为 0.2，以降低输出随机性。

调用失败时，本地模板根据同一结构化输入生成结果。
因此，分类和证据不会受到外部服务状态影响。
LLM 生成内容仍需要教师复核；
系统提供的是建议文本，而不是自动教学决策。
